\documentclass[10pt]{article}
\usepackage[margin=20mm]{geometry}
\usepackage{amsmath,amssymb,fontspec,unicode-math,xurl,hyperref}
\IfFontExistsTF{DejaVu Serif}{\setmainfont{DejaVu Serif}}{\setmainfont{Latin Modern Roman}}
\setmathfont{latinmodern-math.otf}
\setlength{\parindent}{0pt}\setlength{\parskip}{6pt}\sloppy
\begin{document}
ELVIA scientific verification report\par
0055 | I25-40162\par
Position: Societal Impacts Research Requires Benchmarks for Creative Composition Tasks\par
Outcome: Minor; bounded benchmark-coverage, taxonomy-attribution and numerical inaccuracies; released-artifact correspondence remains partly unresolved\par
Batch token: `[internal identifier omitted]`\par
Human scientific decision: Pending\par
\# 1. Identity, coverage and supported outcome\par
The reviewed artifact is the ICML 2025 paper by Judy Hanwen Shen, PMLR 267:82179-82199 [E1]. Printed pages 1-21 equal the 1-based PDF pages. SHA-256: 38d1b06da0a25f3531670e5898e9cf0bf321273aab758f8858a9308a70e5a730. Original bytes match the controller receipt and remain unchanged.\par
All main text and Appendix A were read. The supported overall impact is Minor: several specific descriptions or numbers require correction, while the inspected evidence does not establish a failure of the qualitative task taxonomy or of the central proposal for comprehensive evaluation of quality, novelty and societal harms. The release-to-historical-analysis link remains partly unverified.\par
\# 2. Problem 1 - Prior email/blog evaluation is omitted (Minor)\par
Source claim. Section 3, p3, broadly says creative composition tasks are absent from standardized LLM benchmarks. Figure 1, p4, places Email Writing and Articles and Blogs in a dotted box whose caption describes tasks not measured by prior work. Section 3 on p5 also uses categorical absence language. The figure placement was visually checked.\par
Counterevidence. The official MT-Bench release at commit b494d0c6b4e7935f1764f8439e75da3e66beccc7 is dated 16 June 2023 [E2]. Its 80-question, multi-category evaluation includes the examples below. The released judging templates and README provide single-answer and pairwise scoring, with creativity among the general quality criteria. These are scored evaluation instances predating this manuscript.\par
Constrained everyday email and blog composition falls within the paper’s own Definition 2.1 and examples. Prior inclusion therefore contradicts a literal claim of no measurement for these cells. A factual repair is to add MT-Bench as sparse general-purpose coverage, specify the literature-search scope/cutoff, and distinguish partial task inclusion from dedicated or holistic evaluation.\par
Strongest counterreading and impact. The introduction and conclusion emphasize insufficient coverage; Sections 4 and 6 require quality, novelty and societal-harm assessment together, while Section 7 acknowledges creativity evaluation in arenas. A few MT-Bench items do not establish representative coverage or an adequate joint harm assessment. The verified defect is the binary absence mapping, not a demonstrated failure of that broader insufficiency claim. Minor is supported; the correction is source-verified, but no revised author figure or new evaluation was produced.\par
MT-Bench item | Existing scored task | Relevant manuscript category\par
Q81 | Constrained travel-blog composition | Articles and Blogs\par
Q82 | Professional email requesting feedback | Email Writing\par
Q83 | Product-comparison blog outline | Articles and Blogs\par
Q84 | Persuasive email | Email Writing\par
\# 3. Problem 2 - Professional creativity is misattributed as Big-C (Minor)\par
Section 2.1, p2, attributes creative work by professional artists and writers to Big-C in Kaufman and Beghetto (2009). The cited Four-C model instead distinguishes professional expertise (Pro-c) from eminent contributions (Big-C); see its abstract and pp1-2 [E3]. Profession alone is not the criterion for Big-C.\par
The strongest alternative is that the author intends a practical everyday-versus-professional scope distinction. That is a permissible choice, but it is not the taxonomy attributed to this reference. Correct the terminology and explicitly state whether both professional-level and eminent work are excluded. This is a source-verified textual repair. It does not itself change the operational little-c inclusion rule or establish different sampled themes; the supported impact is a bounded attribution/scope error. The same wording persists in the inspected arXiv v2 passage.\par
\# 4. Problem 3 - The InstructGPT percentage is miscomputed (Minor)\par
Table 1, p5, assigns 57.5\% to InstructGPT generation plus brainstorming. The cited official NeurIPS 2022 supplement, Table 1, physical PDF p1 / printed p16, gives 45.6\% and 11.2\% [E4]. Their sum is 56.8\%, a difference of 0.7 percentage points. Rounding components to one decimal cannot explain 57.5\%. The focal paper’s prose approximation of about 57\% is consistent with 56.8\%.\par
The verified repair is 56.8\% for this stated aggregation, unless a different source population is explicitly documented. This change does not alter the displayed 28\%-62\% range or demonstrate a failure of the prevalence motivation. The upstream categories and sampling frame also differ across studies: broad generation or assisting categories are not identical to the focal definition of creative composition. Table 1 should be read as heterogeneous contextual evidence, not a pooled or population-wide prevalence estimate.\par
\# 5. Problem 4 - Released counts do not fully match the described pipeline (Minor; historical effect Needs Evidence)\par
The two author datasets linked in Appendix A.2 were pinned to LMSYS commit c44782f10f8a329558745fa410686f32d45477b9 and WildChat commit d470c0f576cefc0f29a0b9589f5266e9b205e112 [E5]. Their public Parquet metadata and complete cluster/creative-flag columns were read through verified byte ranges. All 50 clusters per dataset were counted. Cached-byte crosschecks agreed; no author processing code or model was run.\par
Using the manuscript’s rule, a cluster is selected when its count of creative flags divided by its row count exceeds 0.30. The exact integer check is 10 times the creative-flag count > 3 times the row count. No cluster is exactly at 30\%, so changing > to >= cannot resolve the WildChat discrepancy. The released selected pool totals 289,730 rows, consistent with the approximate 289k description, but it contains 42 selected clusters rather than 43.\par
The established issue is a publication/release correspondence discrepancy. Figure 3’s LMSYS count differs by 1,000 rows; Appendix A.2’s WildChat selected-cluster count differs by one. The manuscript’s reported design implies 43*100=4,300 sampled prompt slots, while the pinned labels would imply 4,200 if the same 100-per-cluster rule were applied. Neither number is a verified realized annotation count: exact historical sampled IDs, revisions and coding records were not acquired.\par
Different historical labels or a manual selection override could explain the difference. Consequently no lost theme, fabricated annotation, or changed societal conclusion is inferred. Minor is supported for the observable correspondence defect; its effect on the historical qualitative analysis remains Needs Evidence. Reconcile the data revision and selected cluster/sample manifest before replacing historical totals. The counting procedure is independently checked, but the authors’ release and analysis were not modified.\par
Quantity | Paper | Pinned release\par
LMSYS deduplicated rows, Fig.3 p15 | 423,952 | 422,952\par
LMSYS selected clusters / rows, A.2 p16 | 20 / about154k | 20 /154,173\par
WildChat deduplicated rows, Fig.3 | 375,573 | 375,573\par
WildChat selected clusters / rows, A.2 | 23 / about135k | 22 /135,557\par
\# 6. Checks that preserve the argument and delimit the evidence\par
Sampling and qualitative scope. The two-million figure describes the source-corpus frame. Main Section 2.2 and Appendix A disclose filtering, 50 clusters per dataset, a keyword-based cluster threshold and 100 sampled prompts per selected cluster. Appendix A.3 combines manually identified tasks with Claude summaries. This supports qualitative task discovery, not proof that millions of prompts were manually coded or that a model independently validated the human coding. Section 2.2 limits scope to the included datasets, and p5 footnote 4 explicitly warns that power users distort raw counts.\par
Selection can miss rare or lexically atypical creative uses, and whole selected clusters can contain noncreative prompts. Those limitations constrain representativeness and completeness; they do not negate the existence of observed task types. Appendix A.4 expressly reports residual sexual content after moderation. No hidden claim of perfect removal was inferred.\par
Benchmark and risk scope. The narrow HELM Lite comparison is supported by its dated official description of reference-based generation and missing open-ended instruction following [E7]. BIG-Bench Lite must be distinguished from the full suite. Sections 4-6 generally frame downstream harms as possible or uncertain and propose research directions; no new causal experiment is asserted. This review does not demand such experiments to justify a research-policy position. The adequacy of every proposed metric, every external risk study and every possible benchmark was not established.\par
Version comparison. The retrieved arXiv record’s latest version is v2, dated 24 May 2025 [E6]. Its first page lists Shen and Carlos Guestrin, whereas the PMLR record lists Shen; the PMLR source remains the reviewed artifact. The inspected p2 attribution and pp4,15,16,21 passages retain the relevant content. Selected-page token comparison is only a locator aid because layout reflow changes reading order. No verified repair was found in these checked passages; complete OpenReview history was inaccessible.\par
\# 7. Conclusion and repair/adoption tracking\par
The findings warrant bounded corrections and clearer release provenance. No demonstrated Major failure of the inspected central qualitative or holistic-evaluation claims is asserted.\par
Item | Judgment | Repair verification | Adoption / remaining work\par
Email/blog absence mapping | Minor | Dated scored MT-Bench examples checked | No adoption; revise coverage map/search criteria\par
Big-C attribution | Minor | Four-C primary source checked | No adoption; clarify Pro-c/Big-C scope\par
InstructGPT57.5\% | Minor | Exact sum56.8\% verified | No adoption; identify any alternative denominator\par
Release count correspondence | Minor; historical effect unresolved | Full label-column counts crosschecked | Reconcile revisions and sample manifest\par
Comprehensive evaluation proposal | No demonstrated core failure | Arguments and selected counterevidence checked | Full benchmark/risk validation not performed\par
Human scientific decision | Pending | AI inspection only | Controller acceptance and human judgment pending\par
\# 8. References and reproducible evidence\par
[E1] Shen (2025), PMLR267:82179-82199. https://proceedings.mlr.press/v267/shen25r.html.\par
Zheng et al., MT-Bench / FastChat official release, commit b494d0c6b4e7935f1764f8439e75da3e66beccc7,16June2023. https://github.com/lm-sys/FastChat/tree/b494d0c6b4e7935f1764f8439e75da3e66beccc7/fastchat/llm\_judge. Exact questions, judging prompts, README and commit metadata retained.\par
Kaufman and Beghetto (2009), Beyond Big and Little: The Four C Model of Creativity. Review of General Psychology13(1):1-12. https://doi.org/10.1037/a0013688. Primary abstract retained from https://experts.azregents.edu/en/publications/beyond-big-and-little-the-four-c-model-of-creativity/.\par
Ouyang et al. (2022), Training Language Models to Follow Instructions with Human Feedback. NeurIPS35, supplementary Table1,printedp16. Official main/supplement PDFs and acquisition hashes retained in sources/0055/upstream.\par
Author datasets and immutable revisions are listed in Section5; sources/0055 retains metadata and exact partial Parquet ranges, with offsets and hashes in code\_parquet\_ranges.json. These are partial source payloads, not complete downloaded prompt datasets.\par
Shen and Guestrin, arXiv:2504.06549v2,24May2025. https://arxiv.org/abs/2504.06549v2. [E7] HELM Lite official introduction,19December2023. https://crfm.stanford.edu/2023/12/19/helm-lite.html.\par
Selected upstream LMSYS-Chat-1M ICLR2024 and WildChat arXiv2405.01470v1 sources, release cards and exact read locators are retained in code\_upstream\_audit.json and its acquisition records.\par
Workflow attribution: Paper01 evidence provenance and Paper02 technical review guidance; Kassis et al. (2026), Scientific Agent Skills, https://doi.org/10.48550/arXiv.2609.00065. Workflow source only, not evidence for these findings.\par
\# Appendix. Verification, failures and AI/human roles\par
Reviewer replay for arithmetic: python -X utf8 verification/0055/report\_arithmetic.py. It uses exact decimal and integer calculations over the verified label-count records. Independent data checks and exact range acquisition are documented in code\_audit.md/json; benchmark and inference checks are documented in independent\_claim\_audit.md/json. coverage.json records inspected versus missing work; provenance.json inventories actual files and hashes.\par
The full Four-C PDF download and Sage page returned403 locally; the primary abstract was successfully archived from the author’s institution, and relevant author-copy PDF text was read through web retrieval. The OpenReview forum required browser verification. No inaccessible document is claimed as a locally saved full PDF. Source pp4 and15 and every rendered report page were visually inspected by AI.\par
The primary AI reviewer and independent claim/data reviewers performed these checks. No author code, notebooks, model inference, training, paid API call or new human annotation was executed. No author adoption or human-review declaration is asserted. Self-contained TeX and rendered PDF share the same structured content. A TeX executable was unavailable on PATH; controller compilation, acceptance and human scientific judgment remain pending.\par
\end{document}
